% resumen.tex

\chapter*{Resumen}
\addcontentsline{toc}{chapter}{Resumen}

En Colombia, más de 314.000 familias reciben el servicio de energía eléctrica de manera informal, en barrios donde las conexiones no cumplen ningún estándar técnico ni legal. Estos territorios, conocidos como barrios subnormales (BS), concentran una problemática que va más allá de la infraestructura: la pobreza estructural, la desconfianza hacia los operadores del servicio y la resistencia comunitaria al pago han frustrado durante décadas los intentos de normalización. Esta monografía propone una metodología integral para la normalización de redes eléctricas en BS, que incorpora sistemas de Autogeneración a Pequeña Escala (AGPE) como componente del proceso, dentro del marco del Programa de Normalización de Redes Eléctricas (PRONE), cuya estructura de financiamiento fue modificada en 2024 para habilitar este tipo de proyectos.

El trabajo parte de una caracterización detallada de los BS con énfasis en la región Caribe, donde se concentra el 90\% de los asentamientos certificados del país. Para ello se cruzaron datos del Sistema Único de Información (SUI), la Encuesta de Calidad de Vida del DANE (2024) y los informes de gestión de los planes quinquenales de los Operadores de Red (OR) Air-e y Afinia. Se revisó el marco normativo vigente en materia de normalización de redes y de AGPE, y se sistematizaron tres casos de estudio documentados en la literatura colombiana para identificar qué ha funcionado, qué no y por qué.

Un elemento diferenciador de esta investigación es la consulta directa a 20 profesionales del sector eléctrico y del sector público con experiencia en proyectos de normalización. Sus respuestas revelaron patrones que la literatura técnica no siempre captura: el principal obstáculo para la normalización no es técnico, sino la resistencia de las comunidades al medidor individual y al incremento en la factura. Este hallazgo le da sentido práctico a incorporar el AGPE como estrategia para moderar ese impacto económico y facilitar la aceptación comunitaria del proceso.

Con base en un caso de estudio hipotético construido a partir de las condiciones reales de la región Caribe, se evaluó la viabilidad técnica y financiera de proyectos de normalización con y sin AGPE. Los resultados muestran que el componente de generación renovable puede reducir de forma significativa la factura del usuario normalizado, lo que lo convierte en un argumento técnico y social para mejorar la acogida del proceso en las comunidades.

La propuesta final es una hoja de ruta en cinco fases que orienta a los OR desde el diagnóstico del territorio hasta la sostenibilidad de los activos a largo plazo, definiendo actores responsables, requisitos y entregables para cada etapa. La vigencia 2024, en la que ningún proyecto presentado al PRONE resultó elegible, confirma que esa guía metodológica hace falta.

\bigskip
\noindent\textbf{Palabras clave:} barrios subnormales, normalización de redes eléctricas, autogeneración a pequeña escala, PRONE, metodología, transición energética justa, región Caribe.